\section{Open questions and next-step probes}
\label{sec:open-questions}

Five concrete follow-up probes for anyone extending this replication.

\begin{enumerate}

\item \textbf{Newer QAOA variants (WS-QAOA, RQAOA, ADAPT-QAOA) on the same
instances.}
The replicated paper uses only vanilla QAOA with layered Nelder--Mead.
Since 2020, WS-QAOA (Egger et al.\ warm-start from GW rounding), RQAOA
(Bravyi et al.\ recursive edge-fixing), and ADAPT-QAOA (Zhu et al.\ greedy
operator pool) all report substantial improvements in approximation ratio
at fixed depth. It is an open question whether these variants qualitatively
change the Fig.~7/10--11 success-probability curves, or move any hypothetical
QAOA-vs-classical runtime crossover point.
\emph{Next:} extend \texttt{work/qaoa\_core.py} with the three variants,
re-run C3/C4 sweeps at $p=1..10$, report success\%, $r$, and total function
evaluations vs vanilla. Pure classical, $\sim$2--4 h local wall-clock.

\item \textbf{Modern classical MaxCut baselines vs QAOA.}
Willsch et al.\ compare QAOA to D-Wave and IBM Q, not to state-of-the-art
classical solvers (GW SDP with Burer--Monteiro, branch-and-cut via SCIP,
Tabu / BLS / KHLWG heuristics). The Guerreschi \& Matsuura
``several-hundred-qubit crossover'' claim (the framing referenced in this
slot's brief) depends critically on the classical baseline chosen.
\emph{Next:} add \texttt{work/classical\_baselines.py} wrapping
\texttt{cvxpy} (or SCS) for GW rounding, \texttt{pyscipopt} for exact MILP,
and a numpy Tabu-search. Compare on C4 and on 20 random $G(n,0.5)$ weighted
MaxCut instances at $N \in \{20,40,60,80,100\}$. Emit
ratio-vs-optimum and time-to-target-ratio curves.

\item \textbf{Structured graph families (planar, community, bipartite).}
Willsch et al.\ picked instances with unique ground state and highly
degenerate first-excited states --- a QAOA worst-case-ish scenario. Real
application graphs (VLSI planar, LFR community, structured bipartite) have
exploitable structure that both structured-mixer QAOA and classical
heuristics attack differently.
\emph{Next:} build a small benchmark set (rectangular grid MaxCut,
LFR benchmark graphs, random-regular), run QAOA $p=1..5$ + the classical
baselines above, report whether the $p{=}1$ monotone-improvement trend and
the linear-anneal-init recovery still hold. $\sim$6 h local.

\item \textbf{Noise-model realism for the linear-anneal-init large-$p$ regime.}
Our C5 gap (MaxCut $p{=}10$ success $76.4\%$ vs paper $85.6\%$) is currently
ascribed to optimizer locality on an ideal simulator, but on real hardware
the $p{=}50$ 2-SAT and $p{=}10$ MaxCut runs traverse thousands of two-qubit
gates, so per-gate error $\sim 10^{-3}$ compounds to depth-limiting infidelity.
A noise sweep separates the two effects.
\emph{Next:} add \texttt{work/noise\_channels.py} with Kraus operators for
single-qubit depolarizing ($10^{-4}..10^{-2}$), two-qubit ZZ crosstalk, and a
Gaussian calibration-drift model on gate angles. Re-run C5 under three noise
budgets. Density-matrix sim caps at $\sim 14$ qubits, so full DM for 2-SAT-8A;
Monte-Carlo trajectories for 16-var MaxCut.

\item \textbf{2020--2026 hardware advances and the runtime crossover.}
The paper is a 2019/2020 snapshot; QPU capabilities have since improved
$\sim 10\times$ in qubit count and $\sim 2$--$5\times$ in two-qubit fidelity.
An open question is whether the QAOA runtime-vs-classical crossover point
shifts qualitatively under 2026 hardware parameters (IBM Heron, Quantinuum
H2, surface-code logical qubits), or whether the bottleneck is now algorithmic
(barren plateaus, optimization overhead) rather than gate-level.
\emph{Next:} literature review + analytic re-derivation combining the paper's
runtime accounting with post-2024 gate times and error rates. Pure
literature, no sims. Emit an arXiv-style short note.

\end{enumerate}
